\documentclass[10pt,a4paper]{amsart}
\usepackage[margin=19mm]{geometry}
\usepackage{amsmath,amssymb,mathtools}
\usepackage[scr=rsfs,cal=euler]{mathalfa}
\usepackage{xfrac}
\usepackage[unicode,hidelinks,bookmarks=false]{hyperref}
\newcommand{\ie}{i.e.\ }
\newcommand{\cf}{cf.\ }
\newcommand{\resp}{respectively\ }
\newcommand{\Subsection}{\subsection}
\providecommand{\nobreakdash}{\nobreak\hbox{-}}
\setlength{\emergencystretch}{2em}
\multlinegap0pt
\usepackage{bm}
\usepackage{amssymb}
\usepackage{tikz}
\usepackage{mathtools}
\usepackage[shortlabels]{enumitem}
\usepackage{xcolor}
\usepackage{stackengine}
\usepackage{tcolorbox}
\newtheorem{lemma}{Lemma}
\newtheorem{prop}{Proposition}
\usepackage{cleveref}
\crefname{lemma}{Lemma}{Lemmas}
\crefname{prop}{Proposition}{Propositions}
\newcommand{\Del}{\dot{\Delta}}
\def\N{{\mathbb N}}
\def\R{{\mathbb R}}
\def\Z{{\mathbb Z}}
\newcommand{\abs}[1]{\left\vert#1\right\vert}
\newcommand{\brac}[1]{\ensuremath{\left[ #1 \right]}}
\newcommand{\eps}{\ensuremath{\varepsilon}}
\newcommand{\grad}{\ensuremath{\nabla}}
\newcommand{\iny}{\ensuremath{\infty}}
\newcommand{\norm}[1]{\ensuremath{\left\Vert #1 \right\Vert}}
\newcommand{\smallnorm}[1]{\ensuremath{\Vert #1 \Vert}}
\newcommand{\stardot}{{\mathop{* \cdot}}}
\DeclareMathOperator{\supp}{supp}
\newcommand{\wstar}{\overset{*}{\rightharpoonup}}
\title{Non-decaying weak solutions to the 2D quasi-geostrophic equations}
\author{David M. Ambrose, Ryan Aschoff, Elaine Cozzi,, James P. Kelliher}
\date{Source: arXiv:2609.06836v1 (CC BY 4.0)}
\begin{document}
\maketitle

\noindent\emph{Source and licence.} Non-decaying weak solutions to the 2D quasi-geostrophic equations. Version arXiv:2609.06836v1 (CC BY 4.0), distributed by arXiv under the Creative Commons Attribution 4.0 International licence, http://creativecommons.org/licenses/by/4.0/, as recorded in the arXiv licence metadata for this exact version and shown on the abstract page https://arxiv.org/abs/2609.06836v1. Full licence notice: \texttt{licenses/arxiv-2609.06836-CC-BY-4.0.txt}.\par\smallskip

\noindent\textit{Editorial scope.} Counted unit: the article's prop P:convergence2 (source0.tex lines 2081--2095). The article's complete original proof of it is delivered with it (source0.tex lines 2096--2288, one \texttt{proof} environment of 193 lines). No mathematics is written, repaired or re-derived: the statement and the proof are the article's own lines, in the article's own notation, and no retained line is substituted or re-flowed.  Retained uncounted as the source-local context the proof needs: lines 755--762; lines 866--894; lines 954--985; lines 1661--1663; lines 1701--1704; lines 1759--1767.  Nothing was omitted from any retained span, so the package carries no omission record.  No retained line contains a citation, so the wrapper carries no bibliography and no bibliography entry.  No retained line contains a figure environment, an image-inclusion command or drawing code, so no image is delivered and the package declares no asset dependency.  Editorial formatting only, in addition: an AMS \texttt{amsart} preamble that restates the article's own declarations and macros used by the retained lines, namely the environments \texttt{lemma}, \texttt{prop} on separate counters, together with the standard packages those lines need.  The retained environments print in this document as Lemma 1, Proposition 1 in order of appearance, whereas the article prints the same items with the numbering of its own full document.  The complete native source of the article as distributed in the arXiv e-print for this exact version is bundled unchanged as \texttt{sources/209577/source0.tex}.
\begin{lemma}\label{SnBound}
Assume $f$ belongs to $L^2_{ul}(\R^d)$.  There exists a constant $C>0$ such that for all $n, j\in\Z$,
    \begin{equation*}
        \| S_n f \|_{L^2_{ul}} \leq C\|f \|_{L^2_{ul}},\quad
        \| H_n f \|_{L^2_{ul}} \leq C\|f \|_{L^2_{ul}},\quad
        \| \dot{\Delta}_j f \|_{L^2_{ul}} \leq C\|f \|_{L^2_{ul}}.
    \end{equation*}
\end{lemma}

\begin{lemma}\label{L:WeakStarSectionalConvergence}
    Suppose that  $(f_k)$ is a sequence in $C_b([0, T] \times \R^2)$ with
    \begin{align*}
        f_k \wstar f \text{ in } L^\iny([0, T] \times \R^2)
    \end{align*}
    for some $f \in L^\iny([0, T] \times \R^2)$, with the additional property that for all $h \in C_c^2(\R^2)$ there exists $C_h > 0$ such that
    \begin{align*}
        \abs{\int_{\R^2} h(x) (f_k(t, x) - f_k(s, x)) \, dx}
            \le C_h \abs{t - s}
            \text{ for all } k \in \N, t, s \in [0, T].
    \end{align*}
    Then 
    \begin{align*}
        f_k(t) \wstar f(t)
            \text{ in } L^\iny(\R^2)
            \text{ for all } t \in [0, T].
    \end{align*}
        Moreover, after changing $f(t, \cdot)$ on a set of measure zero in $[0, T]$, for all $s, t \in [0, T]$,
    \begin{align}\label{e:WeakTimeContinuityOff}
        \begin{split}
            &\int_{\R^2} h(x) f(s, x) \, dx
                \to \int_{\R^2} h(x) f(t, x) \, dx
                \text{ for all } h \in L^1(\R^2), \\
            &\abs{\int_{\R^2} h(x) (f(s, x) - f(t, x)) \, dx}
                 \le C_h \abs{t - s}
                 \text{ for all } h \in C_c^2(\R^2).
            \end{split}
    \end{align}
\end{lemma}

\begin{lemma}\label{L:ffInt}
    Fix $N \in \Z$ and for any $\phi \in C_c^\iny(\R^2)$, define the operators $A$, $B$ on scalar functions $f$, $g$ by,
    \begin{align}\label{e:ABOperators}
        \begin{split}
    	A(f, g)
			&:= \int_{\R^2} f(x)
    	    \grad \phi(x) \cdot ((H_N K) * g(x)) \, dx, \\
        B(f, g)
            &= \int_{\R^2} (H_N K) \stardot (f (\grad \phi(\cdot) - \grad \phi(x)))
                    g(x) \, dx.
        \end{split}
    \end{align}
    We have,
	\begin{align}\label{e:ffIntId}
        \begin{split}
            A(f, g)
                &= B(f, g)
                    - A(g, f), \\
            A(f, f)
                &= \frac{1}{2}B(f, f).
            \end{split}
    \end{align}
    Moreover, for a constant $C$ depending on $\phi$, for $f, g \in L^\iny(\R^2)$,
	\begin{align}\label{e:ffBound}
        \begin{split}
        \norm{B(f, g)}_{L^\iny}
            &\le C \norm{f}_{L^\iny} \norm{g}_{L^\iny} 2^{-N}, \\
		\norm{A(f, f)}_{L^\iny}
			&\le C \norm{f}_{L^\iny}^2 2^{-N}.
        \end{split}
	\end{align}
\end{lemma}

\begin{prop}\label{prop:convergence-of-initial-data-u}
Let $\eta\in C^{\infty}_c(\R^2)$.  Then $$ \eta S_k u_0 \rightarrow  \eta u_0 \quad \text{in }L^2(\R^2) \text{ as }k\rightarrow\infty.$$
\end{prop}

\begin{prop}\label{prop:existence-of-limit-of-theta-u}
     Suppose $(\theta_0,u_0) \in L^{\infty}(\R^2)\times L^2_{ul}(\R^2)$ satisfies the homogeneous constitutive law, with div $u_0=0$, and let $(\theta_k,u_k)$ be a solution to $(SQG_{1/k})$ on $[0,T]$ with initial data given by $(\theta_{k,0},u_{k,0})=(S_k\theta_{0}, S_k u_{0})$. Then there exists a pair $(\theta,u)$ such that, up to a subsequence, $$\theta_k\rightarrow\theta \text{ weak-$\star$ in }L^{\infty}([0,T]\times \R^2),$$ and 
     $$u_k \rightarrow u \text{ weak-$\star$ in }L^{\infty}([0,T]; L^2(K)) \text{ for each compact $K$ in $\R^2$}.$$ Further, $$\theta \in L^\infty([0,T]\times\R^2),$$ and $$u \in L^{\infty}([0,T];L^2_{ul}(\R^2)).$$
\end{prop}

    \begin{prop}\label{P:WeakTimeContinuityOfthetaSeq}
        Let $(\theta_k,u_k)$ be as in \cref{prop:existence-of-limit-of-theta-u}.
        For any $h \in C_c^2(\R^2)$ there exists $C_h > 0$ such that
        \begin{align*}
            \abs{\int_{\R^2} h(x) (\theta_k(t, x) - \theta_k(s, x)) \, dx}
                \le C_h \abs{t - s}
                \text{ for all } k \in \N, t, s \in [0, T].
        \end{align*}
    \end{prop}

    \begin{prop}\label{P:convergence2}
    For each fixed $N \in \Z$ and any $\phi \in C_c^\iny([0, T) \times \R^2)$,
    \begin{equation}\label{e:KeyNonlinearConv}
    \begin{split}
    \lim_{k\rightarrow\infty}
    	\int_0^T &\int_{\R^2} \theta_k(t, x)
        \nabla\phi(t,x) \cdot ((H_N K) *
        	(\theta_k(t) - \theta_{k,0}))(x)
            \, dx \, dt \\
        &= \int_0^T\int_{\R^2} \theta(t, x)
            \nabla\phi(t,x) \cdot ((H_N K) * (\theta(t) - \theta_0))(x)
            \, dx \, dt.
    \end{split}
    \end{equation}
    \end{prop} 

    \begin{proof}
    We can write \cref{e:KeyNonlinearConv} in the form
    \begin{align*}
    	\lim_{k \to \iny}
			\brac{A(\theta_k, \theta_k - \theta_{k, 0})
				- A(\theta, \theta - \theta_0)}
			= 0,
    \end{align*}
    where
    \begin{align*}
    	A(f, g)
			&:= \int_0^T \int_{\R^2} f(t, x)
    	    \grad \phi(t,x) \cdot ((H_N K) * g(t))(x)
        	    \, dx \, dt,
    \end{align*}
    which we note is the integrated-in-time form of the operator in \cref{e:ABOperators}. We have,
    \begin{align}\label{e:ABreakdown}
    	\begin{split}
    	A(&\theta_k, \theta_k - \theta_{k, 0})
				- A(\theta, \theta - \theta_0) \\
			&= A(\theta_k, \theta_k - \theta_{k, 0})
				- A(\theta_k, \theta - \theta_0)
				+ A(\theta_k - \theta, \theta - \theta_0) \\
			&= A(\theta_k, \theta_k - \theta)
				+ A(\theta_k, \theta_0 - \theta_{k, 0})
				+ A(\theta_k - \theta, \theta - \theta_0) \\
			&=  A(\theta_k, \theta_0 - \theta_{k, 0})
				+ A(\theta_k - \theta, \theta - \theta_0) \\
			&\qquad\qquad
				+ A(\theta_k - \theta, \theta_k - \theta)
                + A(\theta, \theta_k - \theta).
		\end{split}
    \end{align}
    We bound each of the four terms on the right hand side of 
    \cref{e:ABreakdown} in turn.
    
	\smallskip\noindent$\boldsymbol{A(\theta_k,
		\theta_0 - \theta_{k, 0})}$:
	We have,
	\begin{align*}
		\abs{A(\theta_k, \theta_0 - \theta_{k, 0})}
			&\le \norm{\theta_k}_{L^\iny([0, T] \times \R^2)}
				\norm{\grad \phi}_{L^1(0, T; L^2)}
				\norm{H_N(u_0 - u_{k, 0})}_{L^2(\Sigma)},
	\end{align*}
	where $\Sigma$ is a compact set in $\R^2$ containing the support of $\grad \phi(t)$ for all $t \in [0, T]$. But, for all $k > N$, $H_N (u_0 - u_{k,0}) = H_N(H_k u_0) = H_k u_0 = u_0 - S_k u_0 \to 0$ in $L^2(\Sigma)$ by \cref{prop:convergence-of-initial-data-u}.

    \smallskip\noindent$\boldsymbol{A(\theta_k - \theta,
    	\theta - \theta_0)}$:
    Because $(\theta, u)$ satisfies the homogeneous constitutive law,  $(H_N K) * (\theta(t) - \theta_0)) = H_N (u(t) - u_0) \in L^2_{ul}(\R^2)$, and hence
	\begin{align*}
		\rho(t, x)
			&:= \grad \phi(t,x) \cdot
            	((H_N K) * (\theta(t) - \theta_0))(x)
			\in L^1([0, T] \times \R^2).
	\end{align*}
	\cref{prop:existence-of-limit-of-theta-u} gives
	\begin{align*}
		A(&\theta_k - \theta, \theta - \theta_0)
			= \int_0^T\int_{\R^2} (\theta_k - \theta) (t, x)
				\rho(t, x)
            \, dx \, dt
            \to 0.
	\end{align*}

    \smallskip\noindent$\boldsymbol{A(\theta_k - \theta,
    	\theta_k - \theta)}$:
    Define,
    \begin{align*}
        v_k
            &:= (H_N K) * \theta_k(t)
            = H_N u_k(t), \\
        v &
            := (H_N K) * \theta(t)
            = H_N u(t).
    \end{align*}
    The second expressions for $v_k$ and $v$ hold because $(\theta_k, u_k)$ and $(\theta, u)$ satisfy the homogeneous constitutive law.
	Then, $A(\theta_k - \theta, \theta_k - \theta) \to 0$ can
	be written as
    \begin{align}\label{e:KeyNonFinal}
         \lim_{k \to \iny} &\int_0^T \int_{\R^2}
         	(\theta_k - \theta)(t,x)
            \nabla\phi(t,x) \cdot (v_k - v)(t, x)
                \, dx \, dt
            = 0.
    \end{align}
    
    Define, for $M \ge N$,
    \begin{align}\label{e:HNM}
        H_N^M := \sum_{j = N}^M \Del_j,
    \end{align}
    and let $H_N^M = 0$ for $M < N$. Then $H_N^M K \in L^1(\R^2)$ because each $\smallnorm{\Del_j K}_{L^1} = \smallnorm{\Del_0 K}_{L^1}$.

    Fixing $N$, define
    \begin{align*}
        v_k^M &:= (H_N^M K) * \theta_k(t), \\
        v^M &:= (H_N^M K) * \theta(t).
    \end{align*}
    Using that $H_N^M K \in L^1(\R^2)$ and $\norm{\theta_k(t)} \le \norm{\theta_0}_{L^\iny}$, we know that $v^M \in L^\iny([0, T] \times \R^2)$ and $(v_k^M)_{k = 1}^\iny$ is bounded in $L^\iny([0, T] \times \R^2)$.

    By \cref{prop:existence-of-limit-of-theta-u}, $\theta_k \wstar \theta$ in $L^\iny([0, T] \times \R^2)$ (here, and in what follows, we take subsequences without further comment). By \cref{L:WeakStarSectionalConvergence} with \cref{P:WeakTimeContinuityOfthetaSeq}, then, we know that $\theta_k(t) \wstar \theta(t)$ in $L^\iny(\R^2)$ for all $t \in [0, T]$. Hence, using that $(H_N^M K)(x - \cdot) \in L^1(\R^2)$, for all $(t, x) \in [0, T] \times \R^2$,
    \begin{align*}        
        v_k^M(t, x)
            &= \int_{\R^2} (H_N^M K)(x - y) \theta_k(t, y)
                \, dy \\
            &\to \int_{\R^2} (H_N^M K)(x - y) \theta(t, y)
                \, dy
            = v^M(t, x).
    \end{align*}

    Fixing both $N$ and $M \ge N$, we have
    \begin{align*}
        &\abs{\int_0^T \int_{\R^2} (\theta_k - \theta)(t, x)
            \grad \phi(t,x) \cdot (v_k^M - v^M)(t, x)
            \, dx \, dt } \\
        &\qquad
        \le \norm{\theta_k - \theta}_{L^\iny([0, T] \times \R^2)}
        \norm{\grad \phi}_{L^\iny([0, T] \times \R^2)}
        \smallnorm{v_k^M - v^M}_{L^1(\supp(\grad \phi))} \\
        &\qquad
        \le C \norm{v_k^M - v^M}_{L^1(\supp(\grad \phi))}.
    \end{align*} 
    As we noted, $(v_k^M)_{k = 1}^\iny$ is bounded in $L^\iny([0, T] \times \R^2)$ and $v^M \in L^\iny([0, T] \times \R^2)$, so $v_k^M - v^M$ is bounded above by an $L^1$ function on $\supp(\grad \phi)$: applying the dominated convergence theorem, the pointwise convergence of $v_k^M - v^M$ to 0
    yields
    \begin{align}\label{e:vkMLimitInk}
        \lim_{k \to \iny} \int_0^T &\int_{\R^2}
        	(\theta_k - \theta)(t,x)
            \nabla\phi(t,x) \cdot (v_k^M - v^M)(t, x)
                \, dx \, dt
            = 0.
    \end{align}


	It remains to show that \cref{e:vkMLimitInk} holds in the limit
	as $M \to \iny$---uniformly in $k$---giving \cref{e:KeyNonFinal}.
	Applying \cref{L:ffInt}
	gives
	\begin{align}\label{e:vkMLimitInM}
		\begin{split}
		&\abs{\int_0^T \int_{\R^2}
        	(\theta_k - \theta)(t, x)
            \grad \phi(t,x) \cdot (v_k^M - v^M
            		- (v_k - v))(t, x)
                \, dx \, dt} \\
			&\qquad
			= \abs{\int_0^T \int_{\R^2}
	        	(\theta_k - \theta)(t, x)
    	        \grad  \phi(t,x) \cdot (H_{M + 1} K) *
	        		(\theta_k - \theta))(t, x)
            	    \, dx \, dt} \\
			&\qquad
			\le C \norm{\theta_k - \theta}_
					{L^\iny([0, T] \times \R^2)}^2
				2^{-M}
			\le C 2^{-M},
		\end{split}
	\end{align}
    independent of $k$.
	What we have shown is that
	\begin{align*}
		&\abs{A(\theta_k - \theta, \theta_k - \theta)}
            \le L_{k, M} + C 2^{-M},
	\end{align*}
	where $L_{k, M}$ is the magnitude of the integral in \cref{e:vkMLimitInk}.
    Let $\eps > 0$
	and choose $M$ large enough that $C 2^{-M} < \eps/2$.
	Then choose $n$ large enough that $L_{k, M}  \le \eps/2$ for all
	$k \ge n$ for the given choice of $M$. Then
	$\abs{A(\theta_k - \theta, \theta_k - \theta)} < \eps$ for all $k \ge n$,
	showing that $\abs{A(\theta_k - \theta, \theta_k - \theta)} \to 0$
	as $k \to \iny$.

    \medskip\noindent$\boldsymbol{A(\theta, \theta_k - \theta)}$:
    Recalling that $N$ is fixed, for any $M \ge N$ we can write $H_N = H_N^M + H_M$. Then, applying \cref{L:ffInt}, $A(\theta, \theta_k - \theta) = I_{k, M} + R_M$, where
    \begin{align*}
        I_{k, M}
			:= \int_0^T &\int_{\R^2} ((H_N^M K) * \theta(t, x))  \cdot \grad \phi(t,x)
                     (\theta_k - \theta)(t, x)  \, dx \, dt \\
            &+ \int_0^T \int_{\R^2} ((H_M K) * \theta(t, x))  \cdot \grad \phi(t,x)
                     (\theta_k - \theta)(t, x)  \, dx \, dt
    \end{align*}
    and $\abs{R_M} \le C 2^{-M}$ independently of $k$.
    But, $(H_M K) * \theta = H_M u \in L^{\infty}([0,T];L^2_{ul}(\R^2))$ by \cref{SnBound}, so $((H_M K) * \theta(t, x))  \cdot \grad \phi(t,x) \in L^1([0,T]\times\R^2)$, as also is $((H_N^M K) * \theta(t, x))  \cdot \grad \phi(t,x)$, since $H_M^N K$ is Schwartz-class.
    
    Hence, for any given $M$, $I_{k, M} \to 0$ as $k \to \iny$, though at an unknown rate.
    
    Let $\eps > 0$ and choose $M$ large enough that $\abs{R_M} < \eps/2$.
	Then choose $n$ large enough that $I_{k, M} \le \eps/2$ for all
	$k \ge n$ for the given choice of $M$. Then
	$\abs{A(\theta, \theta_k - \theta)} < \eps$ for all $k \ge n$,
	showing that $\abs{A(\theta, \theta_k - \theta)} \to 0$
	as $k \to \iny$. 
\end{proof}

\end{document}
